\documentclass[12pt]{article}
\usepackage{amsmath,amssymb}
\usepackage{geometry}
\geometry{margin=1in}

\title{CSE 400: Fundamentals of Probability in Computing\\
Tutorial 1 --- Formal Lecture Scribe}
\date{February 13, 2026}

\begin{document}

\maketitle

School of Engineering and Applied Science, Ahmedabad University

\section{Introduction}

This tutorial covers several foundational topics in probability theory relevant to computing systems, including:

\begin{itemize}
\item Multinomial counting
\item Uniform probability models
\item Conditional probability
\item Poisson distribution
\item Gaussian distribution and Q-function
\item Binomial reliability models
\item Discrete distribution functions
\end{itemize}

Each problem illustrates a specific probabilistic model and demonstrates the derivation of the required probability expressions step-by-step.

\section{Multinomial Counting Problem (Food Festival)}

\subsection*{Problem Statement}

Twenty distinct dishes are divided into four platters of five dishes each. The order of dishes within platters and the order of platters themselves are irrelevant.

\begin{enumerate}
\item Number of possible divisions
\item Probability that each platter contains dishes from a single cuisine, given that each cuisine has exactly five dishes.
\end{enumerate}

\subsection*{Solution}

\textbf{(a) Total Number of Arrangements}

There are 20 distinct dishes partitioned into 4 groups of 5 dishes each.

Since:

\begin{itemize}
\item Order within a group does not matter
\item Groups themselves are unlabeled
\end{itemize}

we use the multinomial counting formula.

\[
\text{Total ways} = \frac{20!}{(5!)^4 4!}
\]

\textbf{(b) Probability of Cuisine-Pure Platters}

Each cuisine contains exactly five dishes.

Thus, the only way every platter contains a single cuisine is if each cuisine forms exactly one platter.

Because platters are unlabeled, there is exactly one favorable arrangement.

Therefore

\[
P = \frac{1}{\frac{20!}{(5!)^4 4!}}
\]

which simplifies to

\[
P = \frac{(5!)^4 4!}{20!}
\]

\section{Uniform Distribution Problem (Bus Arrival)}

\subsection*{Problem Statement}

Buses arrive every 15 minutes starting at 7:00 AM. A passenger arrives uniformly at random between 7:00 and 7:30.

Find the probability that the passenger waits:

\begin{enumerate}
\item Less than 5 minutes
\item More than 10 minutes
\end{enumerate}

\subsection*{Solution}

The arrival time is uniformly distributed over a 30-minute interval.

\subsubsection*{(a) Waiting Time $<$ 5 Minutes}

The passenger must arrive within the final 5 minutes before a bus arrives.

Possible intervals:

\begin{itemize}
\item 7:10--7:15
\item 7:25--7:30
\end{itemize}

Total favorable time:

\[
5 + 5 = 10 \text{ minutes}
\]

Total sample space:

\[
30 \text{ minutes}
\]

Therefore

\[
P(\text{wait} < 5) = \frac{10}{30} = \frac{1}{3}
\]

\subsubsection*{(b) Waiting Time $>$ 10 Minutes}

This occurs if the passenger arrives more than 10 minutes before the next bus.

Intervals:

\begin{itemize}
\item 7:00--7:05
\item 7:15--7:20
\end{itemize}

Total favorable time:

\[
10 \text{ minutes}
\]

Thus

\[
P(\text{wait} > 10) = \frac{10}{30} = \frac{1}{3}
\]

\section{Conditional Probability (Employee Attendance)}

Let

\begin{itemize}
\item $L$ = employee arrives late
\item $E$ = employee leaves early
\end{itemize}

Given:

\[
P(L)=0.15, \quad P(E)=0.25, \quad P(L\cap E)=0.08
\]

Find the probability that the employee arrives early given that they leave late.

\subsection*{Solution}

Arriving early is the complement of arriving late:

\[
L^c
\]

Leaving late is the complement of leaving early:

\[
E^c
\]

First compute

\[
P(E^c)=1-P(E)=0.75
\]

Next compute

\[
P(L\cup E)=P(L)+P(E)-P(L\cap E)
\]

\[
=0.15+0.25-0.08
\]

\[
=0.32
\]

Using De Morgan's Law

\[
P(L^c\cap E^c)=1-P(L\cup E)
\]

\[
=1-0.32=0.68
\]

Therefore

\[
P(L^c|E^c)=\frac{P(L^c\cap E^c)}{P(E^c)}
\]

\[
=\frac{0.68}{0.75}
\]

\[
\approx 0.907
\]

\section{Poisson Distribution (Typographical Errors)}

\subsection*{Definition: Poisson Distribution}

A random variable $X$ following a Poisson distribution with mean $\lambda$ has probability mass function

\[
P(X=k)=\frac{e^{-\lambda}\lambda^k}{k!}
\]

\subsection*{Problem}

Expected errors per page:

\[
\lambda=0.2
\]

Find

\begin{enumerate}
\item $P(X=0)$
\item $P(X\ge2)$
\end{enumerate}

\subsection*{(a) Probability of Zero Errors}

\[
P(X=0)=\frac{e^{-0.2}(0.2)^0}{0!}
\]

\[
P(X=0)=e^{-0.2}
\]

\[
P(X=0)=0.8187
\]

\subsection*{(b) Probability of Two or More Errors}

Using the complement rule:

\[
P(X\ge2)=1-P(X=0)-P(X=1)
\]

\[
P(X=1)=\frac{e^{-0.2}(0.2)^1}{1!}
\]

\[
=0.2e^{-0.2}
\]

\[
=0.1637
\]

Thus

\[
P(X\ge2)=1-0.8187-0.1637
\]

\[
=0.0176\approx0.0175
\]

\section{Poisson Distribution (Airplane Crashes)}

Let

\[
X\sim \text{Poisson}(3.5)
\]

\subsection*{(a) At Least Two Accidents}

\[
P(X\ge2)=1-P(0)-P(1)
\]

\[
P(0)=e^{-3.5}=0.0302
\]

\[
P(1)=3.5e^{-3.5}=0.1057
\]

Thus

\[
P(X\ge2)=1-0.0302-0.1057
\]

\[
=0.8641
\]

\subsection*{(b) At Most One Accident}

\[
P(X\le1)=P(0)+P(1)
\]

\[
=0.0302+0.1057
\]

\[
=0.1359
\]

\section{Poisson Distribution (Highway Accidents)}

Let

\[
X\sim \text{Poisson}(3)
\]

\subsection*{(a) Probability $X\ge3$}

\[
P(X\ge3)=1-P(0)-P(1)-P(2)
\]

\[
P(0)=e^{-3}=0.0498
\]

\[
P(1)=3e^{-3}=0.1494
\]

\[
P(2)=\frac{3^2}{2!}e^{-3}=0.2240
\]

Thus

\[
P(X\ge3)=1-0.4232
\]

\[
=0.5768
\]

\subsection*{(b) Conditional Probability}

\[
P(X\ge3|X\ge1)=\frac{P(X\ge3)}{P(X\ge1)}
\]

\[
P(X\ge1)=1-P(0)=0.9502
\]

Thus

\[
P(X\ge3|X\ge1)=\frac{0.5768}{0.9502}
\]

\[
\approx0.607
\]

\section{Gaussian Distribution and Q-Function}

Given PDF

\[
f_X(x)=\frac{1}{\sqrt{8\pi}}\exp\left(-\frac{(x+3)^2}{8}\right)
\]

Comparing with the Gaussian form

\[
\frac{1}{\sqrt{2\pi\sigma^2}}\exp\left(-\frac{(x-m)^2}{2\sigma^2}\right)
\]

we obtain

\[
m=-3,\quad \sigma^2=4,\quad \sigma=2
\]

Relations:

\[
F_X(x)=\Phi\left(\frac{x-m}{\sigma}\right)
\]

\[
P(X>x)=Q\left(\frac{x-m}{\sigma}\right)
\]

Results:

\[
P(X\le0)=\Phi\left(\frac{3}{2}\right)=1-Q\left(\frac{3}{2}\right)
\]

\[
P(X>4)=Q\left(\frac{7}{2}\right)
\]

\[
|X+3|<2 \Rightarrow -5<X<-1
\]

\[
P=\Phi(1)-\Phi(-1)=2\Phi(1)-1
\]

\[
=1-2Q(1)
\]

\[
|X-2|>1
\]

\[
P=F_X(1)+P(X>3)
\]

\[
=1-Q(2)+Q(3)
\]

\section{Jury Decision Problem}

Let

\begin{itemize}
\item $\theta$ = probability a juror makes the correct decision
\item $\alpha$ = probability defendant is guilty
\end{itemize}

If the defendant is innocent, a correct verdict requires at most 7 guilty votes.

\[
\sum_{i=5}^{12}\binom{12}{i}\theta^i(1-\theta)^{12-i}
\]

If the defendant is guilty, a correct verdict requires at least 8 guilty votes.

\[
\sum_{i=8}^{12}\binom{12}{i}\theta^i(1-\theta)^{12-i}
\]

Therefore the total probability of a correct decision is

\[
\alpha\sum_{i=8}^{12}\binom{12}{i}\theta^i(1-\theta)^{12-i}
+
(1-\alpha)\sum_{i=5}^{12}\binom{12}{i}\theta^i(1-\theta)^{12-i}
\]

\section{Poisson PMF Identification}

Given

\[
p(i)=c\frac{\lambda^i}{i!}
\]

Using normalization

\[
\sum_{i=0}^{\infty}p(i)=1
\]

\[
c\sum_{i=0}^{\infty}\frac{\lambda^i}{i!}=1
\]

Since

\[
e^x=\sum_{i=0}^{\infty}\frac{x^i}{i!}
\]

we obtain

\[
ce^\lambda=1
\]

\[
c=e^{-\lambda}
\]

Results

\[
P(X=0)=e^{-\lambda}
\]

\[
P(X>2)=1-e^{-\lambda}-\lambda e^{-\lambda}-\frac{\lambda^2 e^{-\lambda}}{2}
\]

\section{Reliability of Communication Systems}

Model

Each component functions independently with probability $p$.

The system works if at least half of the components function.

\subsection*{(a) Comparison of 5-Component and 3-Component Systems}

For 5 components:

\[
\binom{5}{3}p^3(1-p)^2+\binom{5}{4}p^4(1-p)+p^5
\]

For 3 components:

\[
\binom{3}{2}p^2(1-p)+p^3
\]

The inequality simplifies to

\[
3(p-1)^2(2p-1)>0
\]

Thus

\[
p>\frac{1}{2}
\]

\subsection*{(b) General Result}

For systems with $2k+1$ and $2k-1$ components, the difference in success probabilities is

\[
\binom{2k-1}{k}p^k(1-p)^k(2p-1)
\]

Since the coefficient is positive,

\[
P_{2k+1}>P_{2k-1}
\]

iff

\[
p>\frac{1}{2}
\]

\end{document}